\section{Módulo de Inteligencia Artificial y Validación Experimental}
\label{sec:ml}

La capacidad predictiva del sistema descansa en un módulo de aprendizaje
automático supervisado implementado en Python que integra ingeniería de
características dinámica, generación de etiquetas orientada al movimiento
real del mercado, validación temporal rigurosa y soporte para múltiples
familias de modelos \parencite{goodfellow_deep_2016}. El módulo delega en
cada estrategia la definición de sus propias características, lo que garantiza
que el motor de entrenamiento puede reutilizarse sin modificación mientras las
estrategias evolucionan.

\subsection{Ingeniería de características y generación de etiquetas}

Para construir el conjunto de datos de entrenamiento, el motor invoca el
método abstracto \textit{populate\_indicators()} de la estrategia activa,
que enriquece el marco de datos con los indicadores técnicos propios de esa
estrategia: si la estrategia emplea el índice de fuerza relativa y una media
móvil simple de veinte períodos, el método calcula y añade esas dos columnas;
si otra estrategia solo requiere el índice de fuerza relativa, el marco de
datos resultante tendrá únicamente esa característica adicional. A
continuación, \textit{get\_feature\_columns()} extrae todas las columnas no
pertenecientes al conjunto base precio-volumen (apertura, máximo, mínimo,
cierre, volumen), construyendo así la matriz de características $X$ de forma
dinámica y sin duplicación de código entre estrategias
\parencite{mckinney_python_2022}.

La generación de etiquetas sigue un enfoque prospectivo implementado en el
método \textit{get\_label()}: para cada vela en el instante $t$, la etiqueta
se determina comparando el precio de cierre de la vela siguiente $t+1$ con
el precio de cierre actual. Si el retorno supera un umbral configurable, la
etiqueta es $+1$ (expectativa de subida); si cae por debajo del umbral
negativo, la etiqueta es $-1$ (expectativa de bajada); en caso contrario, la
etiqueta es $0$ (mercado lateral). El umbral por defecto del 0,2\,\%
equilibra la sensibilidad de la señal con la frecuencia de etiquetas activas
\parencite{lopez_de_prado_advances_2018}. Este enfoque es deliberado: si las
etiquetas se derivasen directamente de las condiciones de compra o venta de la
estrategia, el modelo aprendería a memorizar la lógica booleana y no a
predecir el mercado, produciendo métricas artificialmente altas sin
rentabilidad real \parencite{goodfellow_deep_2016}.

El desequilibrio de clases es una consecuencia esperada en datos de
negociación algorítmica, donde la mayoría de los períodos no activan señales
comerciales \parencite{lopez_de_prado_advances_2018}. Este desequilibrio se
aborda en dos dimensiones. Desde el punto de vista de la evaluación, las
métricas se calculan con ponderación por la frecuencia real de cada clase,
evitando que las clases minoritarias sean invisibles en las estadísticas de
rendimiento \parencite{pedregosa_scikit_2011}. Desde el punto de vista del
entrenamiento, todos los modelos ajustan automáticamente el peso de cada clase
mediante \textit{compute\_class\_weight("balanced")} de scikit-learn, de modo
que la clase más escasa contribuye con mayor peso al gradiente de optimización;
cuando el desequilibrio es muy pronunciado, pueden aplicarse adicionalmente
técnicas de remuestreo sintético como SMOTE \parencite{chawla_smote_2002}. En
la práctica, distribuciones muy sesgadas, como el 71\,\% de etiquetas neutras
observado en algunas combinaciones de estrategia y marco temporal, señalan una
incompatibilidad entre la estrategia y el período histórico elegido y deben
motivar una revisión de los umbrales de señal o una ampliación del conjunto de
datos.

\subsection{Preparación del conjunto de datos}

Los datos enriquecidos con indicadores se preparan en tres etapas antes del
entrenamiento. En primer lugar, el período de calentamiento necesario para la
estabilización numérica de los indicadores se determina mediante
\textit{get\_warmup\_period()}, que devuelve el máximo de todos los períodos
configurados en la estrategia; las filas iniciales que no alcanzan ese mínimo
se descartan para evitar valores incompletos que distorsionarían el aprendizaje
\parencite{murphy_technical_1999}. En segundo lugar, se realiza una limpieza
defensiva del conjunto de datos: los valores infinitos se reemplazan por
valores nulos y las filas que los contienen se eliminan, registrando cuántas
filas se descartaron \parencite{mckinney_python_2022}. En tercer lugar, el
conjunto resultante se divide de forma estrictamente cronológica en proporciones
80/20: el ochenta por ciento más antiguo forma el conjunto de entrenamiento y
el veinte por ciento más reciente se reserva como conjunto de evaluación final,
que no participa en ningún momento del proceso de entrenamiento ni de la
búsqueda de hiperparámetros para preservar la causalidad temporal
\parencite{goodfellow_deep_2016,lopez_de_prado_advances_2018}.

Cada modelo entrenado persiste junto con sus metadatos: nombre de la
estrategia, tipo de modelo, lista de características utilizadas, distribución
de etiquetas del conjunto de datos y métricas de clasificación obtenidas. Este
registro permite auditar bajo qué condiciones se obtuvo cada resultado y
detectar el sobreajuste al histórico descrito por
\textcite{lopez_de_prado_advances_2018}. El motor genera adicionalmente un
archivo de valores separados por comas con las métricas de todos los ensayos
de optimización, lo que habilita un análisis posterior de la convergencia del
proceso.

\subsection{Algoritmos soportados y optimización de hiperparámetros}

El módulo admite cuatro familias de modelos seleccionables desde el orquestador
Java sin modificar el código de evaluación. XGBoost \parencite{chen_xgboost_2016}
actúa como modelo de referencia gracias a su capacidad de regularización
explícita, especialmente útil para filtrar indicadores técnicos ruidosos.
LightGBM \parencite{ke_lightgbm_2017} ofrece tiempos de entrenamiento
sustancialmente menores en conjuntos de datos de gran volumen mediante su
estrategia de crecimiento por hojas en lugar de por niveles. Los bosques
aleatorios reducen la varianza del modelo a través del muestreo por remplazo y
aportan una estimación de importancia de características que orienta la
revisión de la ingeniería de características. Las redes neuronales con
TensorFlow\,/\,Keras \parencite{goodfellow_deep_2016} exploran dependencias
complejas en las series de precios mediante arquitecturas configurables de capas
densas con abandono regularizador y detención anticipada. La disponibilidad de
varios modelos facilita la comparación controlada del rendimiento entre familias
de algoritmos bajo condiciones idénticas de entrenamiento y evaluación
\parencite{feurer_hyperparameter_2019}.

Para la búsqueda automática de la configuración óptima, el sistema emplea
Optuna \parencite{feurer_hyperparameter_2019}, una librería de optimización
bayesiana que usa el estimador de Parzen estructurado en árbol para descarta
iterativamente las regiones del espacio de búsqueda con baja probabilidad de
mejora. Un podador por mediana (\textit{MedianPruner}, con cinco ensayos de
arranque y diez pasos de calentamiento) interrumpe de forma anticipada los
ensayos cuyo rendimiento intermedio queda por debajo de la mediana de los
ensayos anteriores, reduciendo el coste computacional total.

La evaluación de cada ensayo sigue un esquema de partición de series
temporales con cinco divisiones por defecto (\textit{TimeSeriesSplit} de
scikit-learn) \parencite{pedregosa_scikit_2011}: en cada partición, el modelo
se entrena sobre el subconjunto cronológicamente anterior y genera predicciones
sobre el subconjunto posterior, de modo que ningún dato futuro contamina la
evaluación. Sobre las predicciones de la partición más reciente ---la que
contiene los datos más próximos al conjunto de evaluación final--- se ejecuta
una simulación histórica para obtener métricas de rendimiento real: tasa de
acierto, factor de ganancia y ratio de Sharpe. La puntuación compuesta que
guía la búsqueda es:
\begin{equation}
    \label{eq:composite_score}
    score = 0{,}25 \cdot F1_{vc} + 0{,}40 \cdot TA + 0{,}25 \cdot FG_{norm} + 0{,}10 \cdot RS_{norm} - \max(0,\, brecha - 0{,}05) \cdot 2
\end{equation}
donde $F1_{vc}$ es el F1 ponderado medio en las particiones de validación,
$TA$ es la tasa de acierto expresada como fracción entre cero y uno,
$FG_{norm} = \min(FG, 5) / 5$ es el factor de ganancia normalizado al
intervalo $[0,1]$, $RS_{norm} = \max(0, \min(RS/2, 1))$ es el ratio de Sharpe
normalizado, y $brecha = acc_{train} - acc_{vc}$ mide el sobreajuste. La
penalización se activa cuando la brecha supera el cinco por ciento,
desincentivando configuraciones que memorizan el conjunto de entrenamiento sin
generalizar \parencite{goodfellow_deep_2016}. La ponderación otorga mayor
peso a la tasa de acierto real (0,40) que a la calidad de clasificación
estadística (0,25), reflejando que el objetivo final es la rentabilidad
operativa y no la optimización de una métrica de aprendizaje aislada
\parencite{lopez_de_prado_advances_2018}.

Una vez completados todos los ensayos, el sistema selecciona los
hiperparámetros del ensayo con mayor puntuación compuesta y entrena con ellos
el modelo definitivo sobre la totalidad del conjunto de entrenamiento,
evaluándolo sobre el conjunto de test reservado para obtener las métricas
finales de exactitud, precisión, exhaustividad y F1. Los espacios de búsqueda
concretos empleados para cada familia de modelos, determinados
experimentalmente durante el desarrollo del proyecto, se recogen en el
Anexo~\ref{anx:hiperparametros}.